
% ============================================================================
\clearpage
\section*{v4.10.0-r25 Source-Only: Two and Four Dixon Copies on the Tensor Shift; the Fourth Copy as the Normal of Three Generations}
\addcontentsline{toc}{section}{v4.10.0-r25 Source-Only: Two and Four Dixon Copies on the Tensor Shift; the Fourth Copy as the Normal of Three Generations}
\label{sec:v4100-r25-dixon-copies}
% ============================================================================

\paragraph{Purpose.}
v4.10.0-r19 found that one generation of \(\C\otimes\Hh\otimes\Oo\) carries no generation qubit: the commutant of gauge and Lorentz group is
\(\C^6\), and only copies give \(M_n(\C)\).  r13 left a three-generation carrier with a neutral Levi--Civita triple open.  The author asked
whether the 64 dimensions of the Dixon algebra suffice, proposed to continue the tensor shift of the atlas (HS(6) is the Dixon carrier with
64 channels) to HS(7) with two and HS(8) with four Dixon algebras, suggested that the fourth algebra store the normals of the other three to
each other, and that the stages partition by halving.  This layer checks these ideas.  Exploration, not blind: the computations are
representation theory and fixed beforehand.  The notes are \texttt{notes/\allowbreak research/\allowbreak v4\_10\_0\_r25/} and
\texttt{notes/\allowbreak research/\allowbreak op3\_dixon\_copies\_2026\_11\_16/}.  Earlier layers are not rewritten.

\noindent\textbf{Source-only status.}
This v4.10.0-r25 layer is source-only.  It does not build a PDF, does not assign a new DOI,
and does not authorize a public release.  The v5.1.0 Version DOI
\href{https://doi.org/10.5281/zenodo.24237112}{10.5281/zenodo.24237112} and the Concept DOI
\href{https://doi.org/10.5281/zenodo.19210728}{10.5281/zenodo.19210728} are unchanged by it.

\subsection*{v4.10.0-r25 results}
\label{subsec:v4100-r25-results}

\paragraph{Theoremlet v4.10.0-r25.1 (commutant on copies).}
One generation (16 Weyl components times 2 spinor components, 32 complex, 64 real) has commutant \(\C^6\) under gauge group and rotations
(as r19).  On \(n\) copies every commutant splits into \(n^2\) blocks, so its dimension is \(6n^2\) (computed for \(n=2\): \(24\); for \(n=4\)
from the block structure: \(96\)), with copy part \(M_n(\C)\).  The gauge group acts trivially on the copy index; every state on it has
definite charge, so the obstacle of r16.1 does not arise there.

\paragraph{Theoremlet v4.10.0-r25.2 (HS(7): two copies give the flat octant).}
The three mutually unbiased states \(+x,+y,+z\) of the copy qubit have overlap \(\tfrac12\), \(K=\tfrac23\), Bargmann phase \(\pi/4\) and
\(\det G=0\): the flat octant of r14, with one massless generation.  Reading: this matches the two generations of route~A in r13 (128
dimensions); that both say the same thing is not shown.

\paragraph{Theoremlet v4.10.0-r25.3 (HS(8): triplet plus singlet; the fourth copy as the normal).}
\begin{enumerate}[leftmargin=2em]
\item[(a)] Four copies are two qubits: \(\C^4=\C^2\otimes\C^2=\mathrm{Sym}^2\oplus\Lambda^2\), a triplet (spin \(1\)) and a singlet (the Levi--Civita
tensor \(\varepsilon_{ab}\)).
\item[(b)] The spin-1 triangle of r15.3 lies wholly in the triplet (singlet part \(10^{-17}\)) and is bent: \(K=\tfrac23\), phase \(0.679042\),
\(\det G=0.050253\).
\item[(c)] The only \(SU(2)\)-equivariant map from three triplet vectors to the singlet is their triple product \(\det V=v_1\cdot(v_2\times v_3)\);
the cross products are the normals.  Computed: \(\det V=-0.224171\,i\) and \(|\det V|^2=\det G=0.050253\), unchanged under a common \(SU(2)\)
rotation; the singlet is invariant and charge-neutral.
\end{enumerate}
\begin{center}
\small
\begin{tabular}{@{}lcc@{}}
\toprule
case & phase \(\varphi\) & \(\det G=|\det V|^2\)\\
\midrule
flat octant (HS(7)) & \(\pi/4\) & \(0\)\\
spin-1 triangle (HS(8)) & \(0.679042\) & \(0.050253\)\\
charged leptons & \(2/3\) & \(0.055706\)\\
\bottomrule
\end{tabular}
\end{center}
The fourth copy is occupied exactly when the three generations are bent, that is when all three masses are positive.  It is the neutral
Levi--Civita carrier sought in r13.6, here on the copy index.

\paragraph{Theoremlet v4.10.0-r25.4 (halving).}
HS(8) \(=\) HS(7) \(\oplus\) HS(7) and HS(7) \(=\) HS(6) \(\oplus\) HS(6): the first copy bit halves HS(8), the second each half; the four boxes
\(|00\rangle,|01\rangle,|10\rangle,|11\rangle\) are the four Dixon algebras.  Triplet and singlet lie across the halving: the singlet
\((|01\rangle-|10\rangle)/\sqrt2\) has weight \(\tfrac12\) in each half.  The carrier has two decompositions, into boxes and into \(3+1\); the
generations are not the boxes but spread over them.

\paragraph{Remark v4.10.0-r25.5 (the boxes carry the mass hierarchy; found afterwards).}
Each of the three generations spreads over the boxes alike: \(|00\rangle\) \(0.7927\), \(|01\rangle\) and \(|10\rangle\) \(0.0976\) each, \(|11\rangle\)
\(0.0120\).  These are the mass-root weights of the spin-1 triangle (\(1:16.24:65.93\)), as the Fourier structure requires.  Against the charged
leptons (\(\tau\) \(0.7931\), \(\mu\) \(0.1934\), \(e\) \(0.01345\)) this is \(0.05\%\), \(0.9\%\) and \(11\%\); the difference is the known residual bend.  Not
blind.

\paragraph{Proof sketch.}
r25.1 by solving the commutation equations (\texttt{dixon\_\allowbreak copies.py}) and the block decomposition.  r25.2 and r25.3 by direct
computation; \(|\det V|^2=\det(V^\dagger V)\).  r25.4 by the product basis of the copy index.  \(\square\)

\paragraph{What this does not give.}
HS(7) and HS(8) are not defined in the atlas; whether the HS(4) gate and the HS(5) debt-free condition admit a seventh and eighth stage is not
checked.  That the generations are read out as coherent spin-1 states with overlap \(\tfrac12\) remains a setting (the core question of
Koide).  The singlet stores the bend but does not fix its value; the residual \(0.012376\) to the lepton phase is not explained.  The frame maps
do not act on the copy index (r19); the reading of r15.5 does not carry over.  No test.

\subsection*{v4.10.0-r25 status in the three frames}
\label{subsec:v4100-r25-frames}

\begin{itemize}[leftmargin=2em]
\item M, T: the copy index is untouched by gauge group and frame maps; triplet, singlet and the triple product are frame-free.
\item P: not evaluated.
\end{itemize}

\subsection*{v4.10.0-r25 register entries}
\label{subsec:v4100-r25-register}

\begin{center}
\small
\begin{tabular}{@{}>{\raggedright\arraybackslash}p{0.34\textwidth}>{\raggedright\arraybackslash}p{0.58\textwidth}@{}}
\toprule
statement & status after v4.10.0-r25\\
\midrule
commutant on \(n\) Dixon copies & \(6n^2\), copy part \(M_n(\C)\), charge-definite (r25.1)\\
HS(7), two copies & copy qubit; three unbiased states form the flat octant (r25.2)\\
HS(8), four copies & triplet \(\oplus\) singlet; spin-1 triangle bent; triple product in the singlet, \(|\det V|^2=\det G\) (r25.3)\\
neutral Levi--Civita carrier (r13.6) & the singlet of four copies, under this reading\\
halving of the tensor shift & HS(8) \(=2\times\)HS(7) \(=4\times\)HS(6); \(3+1\) lies across it (r25.4)\\
HS(7), HS(8) in the atlas & not defined; admissibility open\\
OP3 & open\\
\bottomrule
\end{tabular}
\end{center}

\subsection*{v4.10.0-r25 ledger and guard update}
\label{subsec:v4100-r25-ledger-guard}

\[
R_{\rm op}^{v4100\text{-}r25}=(1,1,0)
\]
for (copy structure recorded; computations checked in the cited script; atlas admissibility and the selection of the states open).

The layer spends no guard:
\[
R_{\rm glob}^{(v4100\text{-}r25,+)}=(0,1,1,1,1),
\qquad
\Gamma_{\rm guard}^{v4100\text{-}r25}=(1,1,1,1,1),
\qquad
\beta_4=1/256\ \text{unspent}.
\]
The layer does not imply exact finite \(C5\) zero closure, exact all-readout terminality, \(C6/T2\), a global photon-clock/RH theorem,
or global all-readout closure.
